\documentclass[letterpaper,11pt]{article}

\usepackage{latexsym}
\usepackage[empty]{fullpage}
\usepackage{titlesec}
\usepackage{marvosym}
\usepackage[usenames,dvipsnames]{color}
\usepackage{verbatim}
\usepackage{enumitem}
\usepackage[hidelinks]{hyperref}
\usepackage{fancyhdr}
\usepackage[english]{babel}
\usepackage{tabularx}
\input{glyphtounicode}

%----------FONT OPTIONS----------
% sans-serif
% \usepackage[sfdefault]{FiraSans}
% \usepackage[sfdefault]{roboto}
% \usepackage[sfdefault]{noto-sans}
\usepackage[default]{sourcesanspro}

% serif
% \usepackage{CormorantGaramond}
% \usepackage{charter}

\pagestyle{fancy}
\fancyhf{} % clear all header and footer fields
\fancyfoot{}
\renewcommand{\headrulewidth}{0pt}
\renewcommand{\footrulewidth}{0pt}

% Adjust margins
\addtolength{\oddsidemargin}{-0.5in}
\addtolength{\evensidemargin}{-0.5in}
\addtolength{\textwidth}{1.0in}
\addtolength{\topmargin}{-.5in}
\addtolength{\textheight}{1.0in}

\urlstyle{same}

\raggedbottom
\raggedright
\setlength{\tabcolsep}{0in}

% Sections formatting
\titleformat{\section}{
  \vspace{5pt}\scshape\raggedright\large
}{}{0em}{}[\color{black}\titlerule \vspace{-5pt}]

% Ensure that generate pdf is machine readable/ATS parsable
\pdfgentounicode=1

%-------------------------
% Custom commands
\newcommand{\resumeItem}[1]{
  \item\small{
    {#1 \vspace{-2pt}}
  }
}

\newcommand{\resumeSubheading}[4]{
  \vspace{5pt}\item
    \begin{tabular*}{0.97\textwidth}[t]{l@{\extracolsep{\fill}}r}
      \textbf{#1} & #2 \\
      \textit{\small#3} & \textit{\small #4} \\
    \end{tabular*}\vspace{-7pt}
}

\newcommand{\resumeSubSubheading}[2]{
  \vspace{-15pt}\item
    \begin{tabular*}{0.97\textwidth}{l@{\extracolsep{\fill}}r}
      \textit{\small#1} & \textit{\small #2} \\
    \end{tabular*}\vspace{-7pt}
}

\newcommand{\resumeProjectHeading}[2]{
    \item
    \begin{tabular*}{0.97\textwidth}{l@{\extracolsep{\fill}}r}
      \small#1 & #2 \\
    \end{tabular*}\vspace{-7pt}
}

\newcommand{\resumeSubItem}[1]{\resumeItem{#1}\vspace{-4pt}}

\renewcommand\labelitemii{$\vcenter{\hbox{\tiny$\bullet$}}$}

\newcommand{\resumeSubHeadingListStart}{\begin{itemize}[leftmargin=0.15in, label={}]}
\newcommand{\resumeSubHeadingListEnd}{\end{itemize}}
\newcommand{\resumeItemListStart}{\begin{itemize}}
\newcommand{\resumeItemListEnd}{\end{itemize}\vspace{-5pt}}

\begin{document}

\begin{center}
    \textbf{\Huge \scshape Chenyu Gao} \\ \vspace{5pt}
    \small
    +1 (667) 910-5300 $|$
    \href{mailto:chenyu.gao@vanderbilt.edu}{chenyu.gao@vanderbilt.edu} $|$
    \href{https://chenyugoal.github.io/}{chenyugoal.github.io} $|$
    \href{https://www.linkedin.com/in/chenyu-gao-4633b31a4/}{LinkedIn} $|$
    \href{https://scholar.google.com/citations?hl=en&user=rKHNyWoAAAAJ&view_op=list_works&sortby=pubdate}{Google Scholar}
\end{center}

\section{Summary}
\small{
Applied ML scientist and final-year Ph.D. in Electrical \& Computer Engineering at Vanderbilt University, expected to graduate in November 2026. Experience at TikTok/ByteDance and insitro spanning LLM post-training, reward modeling for generative video, multimodal representation learning, diffusion models, and biomedical AI. Built large-scale ML/evaluation pipelines and production-quality signals with measurable platform impact, including LLM-driven video understanding for recommendation and creator-content ecosystems. Author of 40+ peer-reviewed publications, including eight first/co-first-author works, with contributions to \textit{Nature} and \textit{IEEE TPAMI}; inventor on one U.S. patent application.
}

\section{Technical Expertise}
\begin{itemize}[leftmargin=0.15in, label={}, itemsep=1pt]
\item \small{
    \textbf{Core AI/ML:} LLM Post-Training/Fine-Tuning (SFT, RL/GRPO), Reward Modeling for Generative Video, Diffusion Models, Vision-Language Models, Multimodal Representation Learning, Self-Supervised Learning, Computer Vision \\
    \textbf{Data Infrastructure \& Evaluation:} Large-Scale Data Organization \& Processing, Human/Model Evaluation \& Benchmarking \\
    \textbf{Engineering:} Python, PyTorch, Hugging Face, Slurm/HPC Workflows, Docker/Singularity, AWS EC2/S3, Git \\
}
\end{itemize}

%-----------EXPERIENCE-----------
\section{Experience}
    \resumeSubHeadingListStart
        \resumeSubheading {Machine Learning Scientist Intern}{San Jose, CA}{TikTok (ByteDance)}{Jan 2026 -- May 2026}
        \resumeItemListStart
            \resumeItem{Designed a reward model for video generation, leveraging GRPO (RL) to improve spatial-temporal consistency and multi-object interaction fidelity while mitigating hallucinations.}
            \resumeItem{Architected a strategy for latent-space reward evaluation using a projection module to align compressed latents with the semantic space of a pre-trained vision foundation model, incorporating object-centric inductive biases.}
            % \resumeItem{Reduced reward-evaluation compute by 75\% by exploiting early DiT semantic-layout emergence, enabling robust video-quality scoring in under 10 denoising steps without pixel decoding.}
            \resumeItem{Reduced computational overhead by cutting required diffusion transformer (DiT) forward steps by 75\%; capitalized on the insight that video semantic layouts emerge during early DiT steps, enabling robust reward evaluation in under 10 steps while also bypassing the pixel-decoding phase.}
            \resumeItem{Developed a lightweight reward model that achieved higher alignment with human raters on target video qualities compared to large Vision-Language Models (VLMs).}
            % \resumeItem{Integrated LLM-driven video understanding into a production Content Ecology pipeline for multidimensional video-quality scoring, driving +4.15\% feature penetration and statistically significant lifts in publish frequency (+0.03\%) and search engagement (+0.019\%).}
            \resumeItem{Integrated LLM-driven video understanding into the Content Ecology pipeline to score multidimensional video quality, inspiring creators and optimizing recommendation algorithms; drove a 4.15\% increase in feature penetration and statistically significant platform-wide lifts in publish frequency (+0.03\%) and search engagement (+0.019\%).}
        \resumeItemListEnd
        \resumeSubheading {Data Science and Machine Learning Intern}{South San Francisco, CA}{insitro}{June 2025 -- Aug 2025}
        \resumeItemListStart
            \resumeItem{Designed and implemented a multimodal framework applying natural language processing (NLP) to encode DNA sequencing data, guiding self-supervised learning (SSL) on images to extract robust features for gene discovery.}
            \resumeItem{Collaborated with research scientists to design model architecture and partnered with software engineers to use and help debug infrastructure pipelines for large-scale computing.}
            \resumeItem{First-authored a technical paper, currently pending internal review for its potential to accelerate drug discovery by identifying new gene targets.}
        \resumeItemListEnd
        
        \resumeSubheading {Graduate Research Assistant}{Nashville, TN}{Vanderbilt University}{July 2022 -- Present}
        \resumeItemListStart
            \resumeItem{Engineered a cascaded diffusion model for a privacy red-team attack, reconstructing high-fidelity 3D facial geometry from defaced MRI data to quantify and expose critical re-identification risks. \href{https://chenyugoal.github.io/research_projects/2024_DREAM/}{[Blog]}}
            \resumeItem{Designed a first-of-its-kind system for brain age estimation from multi-modal MRI to enable early detection of neurodegenerative diseases, resulting in a provisional patent. \href{https://github.com/MASILab/BRAID}{[GitHub]}}
            \resumeItem{Architected a scalable data processing pipeline using Singularity, local and HPC to ingest, harmonize, and quality-assure 100,000 MRI scans from 40+ datasets for the world's largest diffusion and structural MRI database.}
            \resumeItem{Implemented a conditional GAN in PyTorch to synthetically extend the MRI field-of-view, rescuing previously incomplete scans and increasing effective dataset size by an estimated 10-15\%.}
        \resumeItemListEnd

        \resumeSubheading{Graduate Research Assistant}{Baltimore, MD}{Johns Hopkins University}{Dec 2020 -- May 2022}
        \resumeItemListStart
            \resumeItem{Extended the application of a lifelong learning algorithm from vision to speech tasks and validated its omnidirectional knowledge transfer on a spoken digit benchmark, contributing to a publication in the top-tier journal \textit{IEEE TPAMI}. \href{https://github.com/neurodata/ProgLearn}{[GitHub]}}
            \resumeItem{Developed and benchmarked a suite of deep learning and classical computer vision algorithms for MR image analysis, establishing performance baselines that guided subsequent research on MRI defacing.}
        \resumeItemListEnd
      
    \resumeSubHeadingListEnd


%-----------EDUCATION-----------
\section{Education}
  \resumeSubHeadingListStart
    \resumeSubheading
      {Vanderbilt University}{Nashville, TN}
      {Doctor of Philosophy in Electrical \& Computer Engineering}{July 2022 -- Nov 2026 (expected)}
      \resumeItemListStart
      \resumeItem{\textbf{Research Focus:} Applying multimodal representation learning and generative models to solve challenges in medical image analysis and computer vision.}
      \resumeItem{\textbf{Honors:} Graduate School Travel Grant (2023), ECE Day best poster (2025)}
      \resumeItemListEnd
      
    \resumeSubheading
      {Johns Hopkins University}{Baltimore, MD}
      {Master of Science in Biomedical Engineering}{Aug 2020 -- May 2022}
    \resumeSubheading
      {Sun Yat-sen University}{Guangzhou, China}
      {Bachelor of Science in Biomedical Engineering}{Aug 2016 -- June 2020}
  \resumeSubHeadingListEnd


%-----------PUBLICATIONS-----------

\section{Selected Publications \& Patent}
\begin{enumerate}[label=S\arabic*., leftmargin=0.15 in, itemsep=1pt]
    \small{
    \item{
        \textbf{C Gao}*, K Xu*, et al.
        ``Pitfalls of defacing whole-head MRI: re-identification risk with diffusion models and compromised research potential.''
        \textit{Computers in Biology and Medicine}. 2025.
        \textbf{[Co-first author; diffusion-model privacy red-team]}
    }

    \item{
        \textbf{C Gao}, et al.
        ``Brain age identification from diffusion MRI synergistically predicts neurodegenerative disease.''
        \textit{Imaging Neuroscience}. 2025.
        \textbf{[First author; patent-backed multimodal disease prediction]}
    }

    \item{
        \textbf{C Gao}, et al.
        ``Field-of-view extension for brain diffusion MRI via deep generative models.''
        \textit{Journal of Medical Imaging}. 2024.
        \textbf{[First author; generative medical image reconstruction]}
    }

    \item{
        ME Kim, \textbf{C Gao}, et al.
        ``White matter micro-and macrostructure brain charts for the human lifespan.''
        \textit{Nature}. 2026.
        \textbf{[Large-scale neuroimaging collaboration]}
    }

    \item{
        JT Vogelstein, J Dey, \ldots, \textbf{C Gao}, et al.
        ``Simple Lifelong Learning Machines.''
        \textit{IEEE Transactions on Pattern Analysis and Machine Intelligence}. 2025.
        \textbf{[Continual/lifelong learning]}
    }

    \item{
        Y Liu, \textbf{C Gao}, et al.
        ``MetaVoxel: Joint Diffusion Modeling of Imaging and Clinical Metadata.''
        \textit{Medical Imaging with Deep Learning}. 2026.
        \textbf{[Multimodal diffusion modeling]}
    }

    \item{
        \textbf{Patent:} \textbf{C Gao}, BA Landman, ME Kim.
        ``System and Method of Brain Age Identification for Predicting Neuro-Degenerative Disease.''
        U.S. Non-provisional Patent Application No. 19/345,706.
    }
    }
\end{enumerate}
\vspace{-2pt}
\small{\textit{Full publication list:} 40+ peer-reviewed publications, including 8 first/co-first-author works; see Google Scholar.}


%-------------------------------------------
\end{document}
